% Conductors, dielectrics, and circuit components used in chapter 2.
\tikzset{
  metal/.style={pattern=north east lines,pattern color=black!62},
  dielectric/.style={fill=black!7},
}

\newcommand{\EMsecondDiagram}[1]{%
  \adjustbox{max totalsize={.82\linewidth}{.23\textheight}}{%
    \input{tikzfigs/chapter2/fig_#1.tex}%
  }%
}

% Horizontal conductor with its thickness centered at the given height.
\newcommand{\EMplate}[3]{%
  \draw[metal] (#1,{#3-.10}) rectangle (#2,{#3+.10});
}

\newcommand{\EMground}[1]{%
  \begin{scope}[shift={(#1)}]
    \draw (-.25,0)--(.25,0) (-.16,-.10)--(.16,-.10) (-.065,-.20)--(.065,-.20);
  \end{scope}%
}

% Vertical dimension from the bottom coordinate to the top coordinate.
\newcommand{\EMdimY}[4]{%
  \draw[guide] ({#1-.20},#2)--({#1+.20},#2)
    ({#1-.20},#3)--({#1+.20},#3);
  \draw[dim] (#1,#2)--(#1,#3) node[midway,fill=white,inner sep=1pt] {$#4$};
}

% Capacitors with horizontal and vertical terminal directions, respectively.
\newcommand{\EMcap}[3]{%
  \begin{scope}[shift={(#1,#2)}]
    \draw (-.13,-.24)--(-.13,.24) (.13,-.24)--(.13,.24);
    \if\relax\detokenize{#3}\relax\else\node[above] at (0,.25) {$#3$};\fi
  \end{scope}%
}
\newcommand{\EMcapV}[3]{%
  \begin{scope}[shift={(#1,#2)}]
    \draw (-.24,-.13)--(.24,-.13) (-.24,.13)--(.24,.13);
    \node[right] at (.26,0) {$#3$};
  \end{scope}%
}
\newcommand{\EMbattery}[2]{%
  \begin{scope}[shift={(#1,#2)}]
    \draw (-.29,.12)--(.29,.12) (-.17,-.12)--(.17,-.12);
    \node[right,inner sep=1pt] at (.30,.18) {$+$};
    \node[right,inner sep=1pt] at (.30,-.17) {$-$};
  \end{scope}%
}

% Oblique view, with a circular end section projected to an ellipse.
\newcommand{\EMcylinder}[2]{%
  \draw (0,0) ellipse[x radius=.78,y radius=#1];
  \draw (0,#1)--(#2,#1) arc[start angle=90,end angle=-90,x radius=.78,y radius=#1]--(0,-#1);
}
